\section{Problem Tanımı}

Bu çalışmada yıldırım tahmini bir \textbf{ikili sınıflandırma} problemi olarak ele alınmıştır. Amaç, belirli bir günde ve belirli bir istasyon bölgesinde yıldırım olayının gerçekleşip gerçekleşmeyeceğini tahmin etmektir.

\begin{equation}
    y = \begin{cases} 
        1 & \text{yıldırım var} \\
        0 & \text{yıldırım yok}
    \end{cases}
\end{equation}

\section{Veri Birleştirme Yaklaşımı}

Yıldırım verileri koordinat bazlı olduğundan, her yıldırım olayı en yakın meteoroloji istasyonuna eşleştirilmiştir. En yakın istasyon belirlenmesinde Öklid mesafesi kullanılmıştır:

\begin{equation}
    d = \sqrt{(\phi_1 - \phi_2)^2 + (\lambda_1 - \lambda_2)^2}
\end{equation}

burada $\phi$ enlem, $\lambda$ boylam koordinatlarını temsil etmektedir.

\section{Özellik Mühendisliği}

\subsection{Zaman Özellikleri}

Tarih bilgisinden türetilen özellikler:
\begin{itemize}
    \item Ay, gün, yılın günü
    \item Mevsim (döngüsel kodlama)
    \item Haftanın günü
\end{itemize}

Döngüsel özellikler sinüs ve kosinüs dönüşümü ile kodlanmıştır:
\begin{equation}
    x_{sin} = \sin\left(\frac{2\pi \cdot t}{T}\right), \quad
    x_{cos} = \cos\left(\frac{2\pi \cdot t}{T}\right)
\end{equation}

\subsection{Gecikme (Lag) Özellikleri}

Önceki günlerin yıldırım durumu önemli bir gösterge olabilir. Bu nedenle 1, 2, 3 ve 7 günlük gecikme özellikleri oluşturulmuştur:

\begin{equation}
    x_t^{(lag-k)} = x_{t-k}
\end{equation}

\subsection{Hareketli Pencere İstatistikleri}

Son 3, 7 ve 14 günlük dönemler için ortalama ve toplam yıldırım sayıları hesaplanmıştır:

\begin{equation}
    \bar{x}_{t,w} = \frac{1}{w} \sum_{i=1}^{w} x_{t-i}
\end{equation}

\section{Makine Öğrenmesi Algoritmaları}

\subsection{Lojistik Regresyon}

Lojistik regresyon, ikili sınıflandırma için temel bir yöntemdir:

\begin{equation}
    P(y=1|x) = \frac{1}{1 + e^{-(\beta_0 + \beta^T x)}}
\end{equation}

\subsection{Random Forest}

Random Forest, birden fazla karar ağacının oylaması ile çalışır:

\begin{equation}
    \hat{y} = \text{mode}\{h_1(x), h_2(x), ..., h_B(x)\}
\end{equation}

burada $h_b(x)$ her bir karar ağacının tahminidir.

\subsection{XGBoost}

XGBoost, gradient boosting algoritmasının optimize edilmiş versiyonudur:

\begin{equation}
    \hat{y}_i^{(t)} = \hat{y}_i^{(t-1)} + f_t(x_i)
\end{equation}

burada $f_t$ t. iterasyonda eklenen ağaçtır.

\subsection{LightGBM}

LightGBM, büyük veri setleri için optimize edilmiş gradient boosting algoritmasıdır. Leaf-wise büyüme stratejisi kullanır.

\section{Değerlendirme Metrikleri}

Dengesiz veri setleri için uygun metrikler kullanılmıştır:

\begin{itemize}
    \item \textbf{Accuracy (Doğruluk):} $\frac{TP + TN}{TP + TN + FP + FN}$
    \item \textbf{Precision (Kesinlik):} $\frac{TP}{TP + FP}$
    \item \textbf{Recall (Duyarlılık):} $\frac{TP}{TP + FN}$
    \item \textbf{F1-Score:} $2 \cdot \frac{Precision \cdot Recall}{Precision + Recall}$
    \item \textbf{ROC-AUC:} ROC eğrisi altında kalan alan
\end{itemize}

\section{Çapraz Doğrulama}

Model performansının güvenilirliği için 5-katlı stratifiye çapraz doğrulama uygulanmıştır.
